\chapter{Konu Tanımı ve Senaryolar}

\section{Proje Amacı}
Bu proje, gerçek zamanlı sohbet uygulaması için mikroservis mimarisine dayalı bir yazılım sistemi geliştirmeyi amaçlamaktadır. Sistem; birebir ve grup sohbet özelliklerini destekleyen, dosya paylaşımı yapabilen, gerçek zamanlı mesajlaşma sağlayan ölçeklenebilir bir mimari üzerine kurulmuştur.

\section{Proje Kapsamı}
Proje, aşağıdaki temel işlevleri kapsamaktadır:
\begin{itemize}
    \item Kullanıcı yönetimi ve kimlik doğrulama
    \item Birebir ve grup sohbet oluşturma
    \item Gerçek zamanlı mesajlaşma
    \item Mesaj geçmişi saklama ve görüntüleme
    \item Dosya ve medya paylaşımı
    \item Gerçek zamanlı bildirimler
\end{itemize}

Sistem en az 5 mikroservis içermektedir ve bu servisler bağımsız olarak geliştirilebilir, dağıtılabilir ve ölçeklendirilebilir şekilde tasarlanmıştır.

\section{Kullanıcı Senaryoları (User Stories)}

\subsection{Kullanıcı Yönetimi}
\textbf{US-1: Kullanıcı Kaydı}
\begin{itemize}
    \item Bir kullanıcı olarak, sisteme kayıt olabilmeliyim (e-posta, kullanıcı adı, şifre ile).
    \item Sistem, e-posta benzersizliğini kontrol etmeli.
    \item Kayıt sonrası kullanıcı doğrulama e-postası gönderilmeli.
\end{itemize}

\textbf{US-2: Kullanıcı Girişi}
\begin{itemize}
    \item Bir kullanıcı olarak, kayıtlı hesabımla giriş yapabilmeliyim.
    \item Sistem, kimlik bilgilerimi doğrulamalı ve oturum açmalı.
    \item Oturum token'ı döndürülmeli.
\end{itemize}

\textbf{US-3: Profil Yönetimi}
\begin{itemize}
    \item Bir kullanıcı olarak, profil bilgilerimi görüntüleyebilmeliyim.
    \item Profil bilgilerimi (ad, soyad, profil resmi, durum mesajı) güncelleyebilmeliyim.
\end{itemize}

\subsection{Sohbet Yönetimi}
\textbf{US-4: Birebir Sohbet Oluşturma}
\begin{itemize}
    \item Bir kullanıcı olarak, başka bir kullanıcı ile birebir sohbet başlatabilmeliyim.
    \item Sistem, aynı kullanıcılarla mevcut sohbeti varsa yeniden oluşturmamalı.
\end{itemize}

\textbf{US-5: Grup Sohbeti Oluşturma}
\begin{itemize}
    \item Bir kullanıcı olarak, birden fazla kullanıcı ile grup sohbeti oluşturabilmeliyim.
    \item Grup sohbetine katılımcı ekleyip çıkarabilmeliyim.
    \item Grup sohbetine isim verebilmeliyim.
\end{itemize}

\textbf{US-6: Sohbet Listesi Görüntüleme}
\begin{itemize}
    \item Bir kullanıcı olarak, katıldığım tüm sohbetleri görebilmeliyim.
    \item Sohbet listesi son mesaj ve zaman bilgisiyle sıralanmalı.
\end{itemize}

\subsection{Mesajlaşma}
\textbf{US-7: Mesaj Gönderme}
\begin{itemize}
    \item Bir kullanıcı olarak, sohbete metin mesajı gönderebilmeliyim.
    \item Gönderilen mesaj, sohbetteki diğer kullanıcılara gerçek zamanlı olarak iletilebilmelidir.
    \item Mesaj gönderildi/iletildi/okundu durumlarını görebilmeliyim.
\end{itemize}

\textbf{US-8: Mesaj Geçmişi}
\begin{itemize}
    \item Bir kullanıcı olarak, sohbetin mesaj geçmişini görüntüleyebilmeliyim.
    \item Mesaj geçmişi sayfalama (pagination) ile yüklenebilmelidir.
\end{itemize}

\textbf{US-9: Mesaj Silme}
\begin{itemize}
    \item Bir kullanıcı olarak, gönderdiğim mesajları silebilmeliyim.
    \item Silinen mesaj, sohbetteki diğer kullanıcılara "bu mesaj silindi" şeklinde görünmelidir.
\end{itemize}

\subsection{Dosya Paylaşımı}
\textbf{US-10: Dosya Yükleme}
\begin{itemize}
    \item Bir kullanıcı olarak, sohbete dosya (resim, video, doküman) yükleyebilmeliyim.
    \item Sistem, dosya türünü ve boyutunu kontrol etmeli.
    \item Yüklenen dosya için bir URL oluşturulmalı.
\end{itemize}

\textbf{US-11: Dosya İndirme}
\begin{itemize}
    \item Bir kullanıcı olarak, sohbette paylaşılan dosyaları indirebilmeliyim.
    \item Dosya URL'leri güvenli ve zaman sınırlı olmalıdır.
\end{itemize}

\subsection{Bildirimler}
\textbf{US-12: Gerçek Zamanlı Bildirimler}
\begin{itemize}
    \item Bir kullanıcı olarak, yeni mesaj aldığımda gerçek zamanlı bildirim alabilmeliyim.
    \item Bildirimler, çevrimiçi kullanıcılara anında ulaşabilmelidir.
    \item Bildirimler, çevrimdışı kullanıcılar için kuyrukta saklanabilmelidir.
\end{itemize}

\section{Ana Kullanım Akışları}

\subsection{Kullanıcı Kayıt ve Giriş Akışı}
\begin{enumerate}
    \item Kullanıcı kayıt sayfasına girer.
    \item E-posta, kullanıcı adı ve şifre bilgilerini girer.
    \item Sistem kayıt bilgilerini doğrular ve kullanıcıyı oluşturur.
    \item Doğrulama e-postası gönderilir.
    \item Kullanıcı e-postasını doğrular.
    \item Kullanıcı giriş sayfasından e-posta ve şifre ile giriş yapar.
    \item Sistem kimlik bilgilerini doğrular.
    \item Oturum token'ı oluşturulur ve kullanıcıya döndürülür.
    \item Kullanıcı ana sayfaya yönlendirilir.
\end{enumerate}

\subsection{Mesaj Gönderme Akışı}
\begin{enumerate}
    \item Kullanıcı bir sohbet seçer.
    \item Mesaj yazma alanına metin girer.
    \item Mesajı gönder butonuna tıklar.
    \item İstemci, mesajı Message Service'e gönderir.
    \item Message Service mesajı doğrular ve veritabanına kaydeder.
    \item Message Service, Notification Service'e bildirim gönderir.
    \item Notification Service, sohbetteki diğer kullanıcılara WebSocket üzerinden bildirim gönderir.
    \item Çevrimiçi kullanıcılar mesajı gerçek zamanlı olarak görür.
    \item Çevrimdışı kullanıcılar için bildirim kuyrukta saklanır.
\end{enumerate}

\subsection{Dosya Paylaşımı Akışı}
\begin{enumerate}
    \item Kullanıcı sohbet içinde dosya ekle butonuna tıklar.
    \item Dosya seçer (resim, video, doküman).
    \item İstemci, dosyayı File Service'e yükler.
    \item File Service dosyayı doğrular (tür, boyut kontrolü).
    \item File Service dosyayı depolama sistemine kaydeder.
    \item File Service dosya URL'ini döndürür.
    \item İstemci, dosya URL'ini içeren bir mesaj oluşturur.
    \item Mesaj normal mesaj gönderme akışı ile gönderilir.
    \item Alıcılar dosyayı görüntüleyebilir veya indirebilir.
\end{enumerate}

\subsection{Grup Sohbeti Oluşturma Akışı}
\begin{enumerate}
    \item Kullanıcı "Yeni Grup" butonuna tıklar.
    \item Grup adını girer ve katılımcıları seçer.
    \item Chat Service'e grup oluşturma isteği gönderilir.
    \item Chat Service grubu oluşturur ve veritabanına kaydeder.
    \item Chat Service, seçilen kullanıcılara grup davet bildirimi gönderir.
    \item Bildirimler Notification Service üzerinden gönderilir.
    \item Grup sohbeti oluşturulur ve kullanıcılar sohbet listesinde görür.
\end{enumerate}
